\chapter{\fem}
\label{chp:fem}

In questo capitolo viene descritta l'azienda \textbf{\femsrl} in modo da contestualizzare i bisogni che hanno portato allo sviluppo della piattaforma web dedicata.

In particolare nelle sezioni~\ref{sec:storia-fem} e~\ref{sec:richiesta} sono illustrate prima nascita ed espansione della spin-off, poi bisogni, richieste e necessità legate alla piattaforma sviluppata.

\section{La storia di \fem}
\label{sec:storia-fem}

\begin{wrapfigure}{r}{0.32\textwidth}
  \begin{center}
    \includegraphics[width=0.3\textwidth]{images/logo-fem}
  \end{center}
\end{wrapfigure}

\textbf{\femsrl} è uno spin-off del Dipartimento di Biotecnologie e Bioscienze dell'{\unimib}, nato con l'intenzione di creare prodotti e servizi per il largo pubblico finalizzati alla conoscenza e tutela della biodiversità. La mission è supportare i consumatori nelle scelte, rendendoli consapevoli sulla qualità delle risorse ambientali, e in questo modo fornendo gli strumenti necessari a migliorare il loro stile di vita, tutelando l'ambiente \cite{fem2ambiente}.

Ad inizio 2007 é nato \textbf{\zpl} dall'incontro di Massimo Labra e Maurizio Casiraghi, due ricercatori del Dipartimento di Biotecnologie e Bioscienze dell'Università degli Studi di Milano Bicocca che si occupano rispettivamente di tematiche botaniche e zoologiche. Lo {\zpl} é un laboratorio di ricerca di zoologia, botanica e microbiologia che coniuga ricerca di base e applicata con progetti che prevedono l'utilizzo di un approccio molecolare \cite{zooplantlab}.

Nel gennaio 2010 grazie al contributo e i risultati della ricerca di {\zpl} viene fondato {\femsrl} dai quattro soci fondatori: Dott. De Mattia Fabrizio, Dott. Ferri Emanuele, Dott. Labra Massimo, Dott. Casiraghi Maurizio.

Food, Environment \& ManageMent (FEM2): alimentazione, ambiente e gestione razionale sono alcuni degli aspetti ai quali {\fem} dedica particolare attenzione, con l'ottica di fornire informazioni e strumenti per un utilizzo più consapevole delle risorse, in sintonia con il pianeta, utilizzando tecnologie e conoscenze derivanti dalla ricerca scientifica.

{\fem} dispone di moderni laboratori ospitati presso l’Università, nei quali vengono sviluppati e testati i nuovi prodotti, eseguite analisi su matrici ambientali (es: acqua, aria o alimenti), si svolgono analisi del DNA e vengono messe a punto metodiche innovative di caratterizzazione molecolare. Grazie ad essi oggi {\fem}, pur mantenendo le sue solide radici universitarie ed investendo nella ricerca, si è affermata anche come società commerciale e propone al mercato nazionale ed internazionale prodotti e servizi all'avanguardia nei settori dell'ambiente, del food e della diagnostica molecolare avanzata.

Negli ultimi anni è diventato leader di mercato nella \emph{diagnostica molecolare di avifauna} tramite PCR (analisi del DNA), ed é nata la necessita di sviluppare una piattaforma adatta per gestire tutte le fasi di analisi e vendita dei servizi.

\section{La richiesta}
\label{sec:richiesta}

La crescita di {\fem} sul mercato ha portato alla creazione del sito dedicato \url{www.fem2ambiente.com} basato su \emph{Joomla!}, un \emph{CMS} (content management system) molto diffuso, utile per la gestione dei contenuti del sito web senza la necessità di avere conoscenze tecniche \cite{joomla}. Ad esso é stata aggiunta l'estensione \emph{VirtueMart}, una soluzione open-source per la gestione dell'e-commerce che ad oggi permette l'acquisto di prodotti per l'analisi di acqua e aria rivolti a privati, condomini, imprese e per l'educazione nelle scuole, oltre ad eco-prodotti per la casa e per la cura degli animali \cite{virtuemart}.

Quando {\fem} ha cominciato l'analisi su avifauna l'utilizzo di fogli elettronici Excel é sembrata la soluzione migliore, ma col crescere della clientela e del numero di ordini si é dovuto pensare ad una alternativa. 

É nata così la richiesta di una piattaforma online dedicata in grado di offrire servizi di \emph{sessaggio}, \emph{diagnosi patologie} e \emph{Dna barcoding}. In particolare in grado di gestire la vendita specifica delle analisi e servizi offerti (non realizzabile attraverso un semplice tool e-commerce già esistente), e il controllo del flusso ordini integrato con i procedimenti in laboratorio.

Sono stati sviluppati di conseguenza
\begin{itemize}
\item un sito dedicato all'esposizione dei servizi offerti, il \emph{Portale della Diagnostica Molecolare dedicato all'Avifauna} (capitolo~\ref{sec:wp})
\item una piattaforma lato server per il controllo completo del flusso ordini (capitolo~\ref{sec:server})
\item un interfaccia per il flusso cliente durante la creazione e monitoraggio degli ordini e un pannello admin per la gestione completa di ogni parte della piattaforma (capitolo~\ref{sec:client})
\end{itemize}
